\documentclass[a4paper,12pt]{article}
\usepackage[russian]{babel}

% Самые необходимые пакеты
\usepackage[T2A]{fontenc}
\usepackage[utf8]{inputenc}
\usepackage[russian]{babel}
\usepackage{geometry}
\usepackage{amsmath}
\usepackage{graphicx}
\usepackage{float} 
\usepackage{xcolor}
\usepackage{listings}
\usepackage{hyperref}
\usepackage{array}
\usepackage{caption}
\usepackage{booktabs}
\usepackage{amssymb}
\usepackage{subcaption}
\usepackage{caption}
\captionsetup[figure]{name=Рисунок}
\captionsetup[table]{name=Таблица}
\usepackage{pgfplots}
\pgfplotsset{compat=1.18}
\usetikzlibrary{arrows.meta}
\usepackage{capt-of}
\usepackage{multirow}
\usepackage[siunitx]{circuitikz}
\usepackage[bottom]{footmisc}
\usepackage{tcolorbox}
\usetikzlibrary{arrows.meta, patterns}

\usepackage[x11names, svgnames]{xcolor}
\definecolor{Lavender}{RGB}{160,110,230}
\definecolor{Peach}{RGB}{255,120,60}
\definecolor{Apricot}{RGB}{255,100,30}
\definecolor{Salmon}{RGB}{255,60,40}
\definecolor{CornflowerBlue}{RGB}{30,80,255}
\definecolor{SlateGray}{RGB}{80,90,110}


\lstset{
	language=Matlab,
	basicstyle=\ttfamily\small,
	keywordstyle=\color{blue}\bfseries,
	commentstyle=\color{gray}\itshape,
	stringstyle=\color{red},
	numberstyle=\tiny\color{gray},
	numbers=left,
	stepnumber=1,
	numbersep=8pt,
	backgroundcolor=\color{white},
	frame=single,
	breaklines=true,
	breakatwhitespace=true,
	tabsize=4,
	showspaces=false,
	showstringspaces=false,
	captionpos=b,
	escapeinside={\%*}{*)},
}
\geometry{left=1.5cm, right=2cm, bottom=2cm, top=2cm}

\begin{document}
	
	% Титульный лист
	\begin{titlepage}
		\centering
		
		% Министерство
		\vspace*{0.5cm}
		{\small
			Министерство науки и высшего образования Российской Федерации\\[2pt]
			Федеральное государственное автономное образовательное учреждение\\
			высшего образования\\[2pt]
			\textbf{«Национальный исследовательский университет ИТМО»}\\
		}
		
		\vspace{0.8cm}
		\hrule height 0.8pt
		\vspace{0.3cm}
		
		{\large Факультет систем управления и робототехники}
		
		\vspace{0.3cm}
		\hrule height 0.4pt
		
		\vspace{2.5cm}
		
		% Основной заголовок
		{\Huge\bfseries ОТЧЁТ}\\[0.4cm]
		{\LARGE\bfseries о курсовой работе}\\[0.8cm]
		
		% Декоративная рамка вокруг темы
		\begin{tcolorbox}[
			colback=gray!4,
			colframe=gray!10,
			boxrule=0.8pt,
			arc=4pt,
			left=12pt, right=12pt, top=8pt, bottom=8pt
			]
			\centering
			{\large По дисциплине:\\[4pt]}
			{\large\bfseries «Теория автоматического управления»}\\[10pt]
			{\large Тема:\\[4pt]}
			{\large\bfseries «Управление нелинейным объектом методами линейных систем»} \\[10pt]
			{\large Поток: 1.1.3}
		\end{tcolorbox}
		
		\vspace{3cm}
		
		% Студент / Преподаватель
		\begin{minipage}{0.45\textwidth}
			\raggedright
			{\large\bfseries Студент:}\\[4pt]
			Охрименко Е.\,Д.\\
		\end{minipage}
		\hfill
		\begin{minipage}{0.45\textwidth}
			\raggedleft
			{\large\bfseries Преподаватель:}\\[4pt]
			Пашенко А.\,В.\\
		\end{minipage}
		
		\vfill
		
		
		\vspace{0.4cm}
		
		{\large Санкт-Петербург, 2026}
		
	\end{titlepage}
	
	\newpage
	\tableofcontents
	
	
	\newpage
	\section{Построение математической модели объекта}
	
	\subsection{Вывод уравнения}


\begin{figure}[h]
\begin{center}
	\begin{tikzpicture}[
		scale=1.5,
		>={Stealth[scale=1.5]}
		]
		
		% Пол
		\draw[->, thick] (-3,-0.2) -- (4,-0.2);
		
		% Тележка
		\draw[thick, fill=Lavender!40] (-1.2,0) rectangle (1.2,0.6);
		
		% Колеса
		\filldraw[SlateGray!80] (-0.7,-0.035) circle (0.15);
		\filldraw[SlateGray!80] (0.7,-0.035) circle (0.15);

		
		% Маятник
		\draw[thick] (0,0.6) -- (-1.5,2.5);
		\filldraw[SlateGray!80] (-1.5,2.5) circle (0.06);
		
		% Ось
		\draw[dashed, very thick] (0,2.7) -- (0,-0.6);
		\draw[thick] (-2,-0.1) -- (-2,-0.3);
		
				
		% Шарнир
		\filldraw[SlateGray!80] (0,0.6) circle (0.06);
		
		% Углы
		\draw[->, thick, Salmon, -{Stealth[length=1.5mm, width=1.2mm]}] (0.,1.22) arc (50:145:0.3);
		\draw[->, thick, Salmon, -{Stealth[length=1.5mm, width=1.2mm]}] (-1.2,2.11) arc (120:160:0.5);
		
		% Сила
		\draw[->, Salmon, thick] (1.215,0.3) -- (2.5,0.3);
		
		% Подписи
		\node at (4,-0.5) {$x$};
		\node[Salmon] at (1.8,0.5) {$u$};
		\node[Salmon] at (-0.2,1.5) {$\varphi$};
		\node[Salmon] at (-1.48,2.1) {$f$};
		\node at (-1.1,1.4) {$L, m_1$};
		\node at (-1.97,2.6) {$m_2$};
		\node at (0.3,-0.5) {$\alpha$};
		\node at (-1.8,-0.5) {$0$};
		\node at (-0.4,0.3) {$M$};
		
		
	\end{tikzpicture}
\end{center}
\caption{Перевернутый маятник на тележке}
\end{figure}
		
	Запишем уравнения движения системы, применяя уравнения Лагранжа второго рода
	для обобщённых координат $q_1 = a$, $q_2 = \varphi$.
	Обозначим $m = m_1 + m_2$, $\beta = \left(\tfrac{m_1}{2} + m_2\right)l$,
	$J = \tfrac{1}{3}m_1 l^2 + m_2 l^2$.
	
	\begin{equation}
		\begin{cases}
			\dot{x}_1 = x_2,\\[4pt]
			\dot{x}_2 = \dfrac{J\bigl(u + \beta x_4^2 \sin x_3\bigr)
				- \beta\cos x_3\bigl(f + \beta\tfrac{g}{l}\sin x_3\bigr)}
			{(M+m)J - \beta^2\cos^2 x_3},\\[10pt]
			\dot{x}_3 = x_4,\\[4pt]
			\dot{x}_4 = \dfrac{(M+m)\bigl(f + \beta\tfrac{g}{l}\sin x_3\bigr)
				- \beta\cos x_3\bigl(u + \beta x_4^2 \sin x_3\bigr)}
			{(M+m)J - \beta^2\cos^2 x_3},\\[10pt]
			y_1 = x_1, \quad y_2 = x_3.
		\end{cases}
	\end{equation}
	
	где $x_1 = a,\ x_2 = \dot{a},\ x_3 = \varphi,\ x_4 = \dot{\varphi}$ —
	вектор состояния, $u$ — управляющая сила, $f$ — возмущающий момент.	
		
	\subsection{Точки равновесия}
	
	При $u = 0$, $f = 0$, $\dot{x}_i = 0$ из уравнений движения следует:
	\[
	\sin\varphi = 0 \implies \varphi \in \{0,\ \pi\}.
	\]
	Координата тележки $a$ произвольна. Получаем два семейства равновесий:
	$P_1$ ($\varphi = 0$, маятник вверх) — \textbf{неустойчивое},
	$P_2$ ($\varphi = \pi$, маятник вниз) — \textbf{устойчивое}.
	Задача проекта — стабилизация в точке $P_1$.
	
	\subsection{Линеаризация}
	
	Линеаризуем систему около $P_1$. При малых отклонениях:
	$\sin\varphi \approx \varphi$, $\cos\varphi \approx 1$, $\dot{\varphi}^2 \approx 0$.
	Получаем модель:
	\[
	\dot{x} = Ax + Bu + Df, \quad y = Cx,
	\]
	где $x = (a,\ \dot{a},\ \varphi,\ \dot{\varphi})^\top$, и
	
	\[
	A = \begin{pmatrix}
		0 & 1 & 0 & 0 \\
		0 & 0 & -\dfrac{\beta g}{M J} & 0 \\[6pt]
		0 & 0 & 0 & 1 \\
		0 & 0 & \dfrac{(M+m)g\beta}{MJ} & 0
	\end{pmatrix}, \quad
	B = \begin{pmatrix} 0 \\ \dfrac{J}{MJ} \\[6pt] 0 \\ -\dfrac{\beta}{MJ} \end{pmatrix},
	\]
	
	\[
	D = \begin{pmatrix} 0 \\ -\dfrac{\beta}{MJ} \\[6pt] 0 \\ \dfrac{M+m}{MJ} \end{pmatrix}, \quad
	C = \begin{pmatrix} 1 & 0 & 0 & 0 \\ 0 & 0 & 1 & 0 \end{pmatrix},
	\]
	где $M_\text{eff} = (M+m)J - \beta^2$ вместо $MJ$ в знаменателях.	
		
	\subsection{Выбор параметров}
	
	По результатам генерации получены следующие параметры системы:
	\[
	M = 658.36 \text{ кг}, \quad
	m_1 = 15.66 \text{ кг}, \quad
	m_2 = 14.80 \text{ кг}, \quad
	l = 1.77 \text{ м}, \quad
	g = 9.81 \text{ м/с}^2.
	\]
	Вспомогательные параметры:
	\[
	m = m_1 + m_2 = 30.46 \text{ кг}, \quad
	\beta = \left(\tfrac{m_1}{2} + m_2\right)l = 39.94, \quad
	J = \tfrac{1}{3}m_1 l^2 + m_2 l^2 = 62.47.
	\]
	Матрицы линеаризованной модели:
	\[
	A = \begin{pmatrix}
		0 & 1 & 0 & 0 \\
		0 & 0 & -0.0095 & 0 \\
		0 & 0 & 0 & 1 \\
		0 & 0 & 6.5197 & 0
	\end{pmatrix}, \quad
	B = \begin{pmatrix} 0 \\ 0.0015 \\ 0 \\ -0.0010 \end{pmatrix},
	\]
	\[
	D = \begin{pmatrix} 0 \\ -0.0010 \\ 0 \\ 0.0166 \end{pmatrix}, \quad
	C = \begin{pmatrix} 1 & 0 & 0 & 0 \\ 0 & 0 & 1 & 0 \end{pmatrix}.
	\]	
		
		
		
		
		
		
		
		
		
		
		
		
		
		
		
		
		
		
		
		
	\end{document}
	
	
	
	
	
	
	
	
	
	
	
	
	
	
	
	
	
	
	
	
	
	
	
	
	
	
	
	
	
	
	
	
	
	
	
	
	
	
	
	
	
	
	
	
	
	
	
	
	
	
	
	
	
	
	
	
	
	
	
\end{document}